%% capability-ladder.tex — the four-rung capability ladder (Chapter 2). House conceptual style.
\input{figures/concept-style}
\begin{tikzpicture}[
    rung/.style={conceptbox, minimum width=3cm, font=\small\bfseries},
    lbl/.style={font=\scriptsize, align=left, text width=6.2cm},
]
\node[rung, fill=gray!15]   (op)  at (0,0)    {Operate};
\node[rung, fill=blue!15]   (bld) at (0,1.25) {Build};
\node[rung, fill=orange!15] (orc) at (0,2.5)  {Orchestrate};
\node[rung, fill=green!15]  (evl) at (0,3.75) {Evolve};
\draw[conceptflow] (op)  -- (bld);
\draw[conceptflow] (bld) -- (orc);
\draw[conceptflow] (orc) -- (evl);
\node[lbl, right=0.6cm of op]  {Terminal, shell, version control, model API};
\node[lbl, right=0.6cm of bld] {Tool definitions, the agent loop, an evaluation harness};
\node[lbl, right=0.6cm of orc] {Multi-agent workflows, retrieval, reliability};
\node[lbl, right=0.6cm of evl] {generate $\rightarrow$ evaluate $\rightarrow$ select $\rightarrow$ mutate};
\end{tikzpicture}
